% !TeX root = ../papers/nc-p3.tex
\section{量子四面の単純値・単一重複値を持つ全退化ファイバー}
\label{sec:quantum-face-resonance}
$q_{ab}$ は $x_bx_a=q_{ab}x_ax_b$ の規約とし、向き付き四面の値を
\[
 (h_0,h_1,h_2,h_3)=(T_{123},T_{023}^{-1},T_{013},T_{012}^{-1})
\]
とする。積は1で、奇置換は添字の置換に全値の逆数化を伴う。
\begin{theorem}\label{thm:quantum-face-resonance}
量子点を軌道閉包に持つ窓の全添字保持同型類について、次が成立する。
\begin{enumerate}
\item 四つの $h$ が相異なれば、この量子点自身だけである。
\item ちょうど一対 $h_k=h_l=c\ne1$ が等しく、残る $h_i,h_j$ は互いにも $c$ にも異なるとする。
このとき背景の六関係のうち $kl$ の一関係だけを変える、下記の三係数表示が全類を尽くす。
同型は各非零台と、三係数全てが非零なら $C_0C_2/C_1^2$ で決まる。
$c=-1$ の場合に限り台は $s\leftrightarrow2-s$ でも同一視する。
零係数を含む全ての類で最小周期は1、標準Koszul左右AS正則である。
\end{enumerate}
異なる背景間の重複は定理\ref{thm:quantum-indexed-all}の背景同型だけを考慮すればよい。
\end{theorem}

\begin{proof}[四面の値で決まる横断指標]
各行積を1にした量子座標を用いる。四位置のある量子窓に対する一次変形で、
関係を順序二次単項式へ書き直す。関係二文字の多重次数からの差を $\alpha$ とする。
$\sum\alpha_j=0$ で、負成分は $-1$ だけ、負の個数は高々二つである。
負の添字集合を $N$ とする。量子Laurent環で $x_j x^\alpha=\chi_j x^\alpha x_j$ と書く。
位置方向を $\eta^n$ ($\eta^4=1$) に分ける。$a,b$ 関係の一次補正を
$\eta^n c_{ab}x^\alpha x_bx_a$ の順で書けば、$c_{ab}$ は $N\subset\{a,b\}$ でのみ許される。
三文字重なりの線型条件は、$v_j=1-\eta\chi_j$ として $v\wedge c=0$ である。
実際、三文字を逆順から整列する二通りの差で、補正対の外の文字 $j$ を右に置いた項の係数は1、
左に置いた項は位置が一つずれて $x^\alpha$ を越えるため $\eta\chi_j$ になる。
背景の交換係数は共通の順序語で相殺する。
一次の文字変更も同じ計算で $c\mapsto c+v\wedge b$、$b_a$ は $N\subset\{a\}$ のときだけ許される。
この線型条件は少なくとも必要である。四窓の階数を保つ量子点の一次変形は各経路で自由なので
三文字で非零の重なりを生じない。関係の変化を基底変更で消せれば、
四次双対縮約直線の一意性からテンソル側にも別の射影変化は残らない。

負添字の外積を先に固定すると、残る条件と同値は相補添字上の $v$ との外積複体になる。
一成分でも非零ならその双対座標による内部積で可縮である。
従って $\alpha\ne0$ の横断指標が出るには $j\notin N$ の $\chi_j$ が全て同じでなければならない。
$N=\{b\}$ では $\alpha=e_a-e_b$ で、これは $a,b$ を含む二つの三角積が
ともに1となる条件である。定理の二場合では面値1が高々一つなので起こらない。
$N=\{k,l\}$ では $\alpha_s=se_i+(2-s)e_j-e_k-e_l$、$s=0,1,2$ である。
$\chi_i/\chi_j$ は $T_{ijk}T_{ijl}$ またはその逆なので、$\chi_i=\chi_j$ は $h_k=h_l$ に等しい。
このとき $\prod\chi_a=1$ と $\prod\chi_a^{\alpha_{s,a}}=1$ により $\chi_i^4=1$ である。
許されるモードは $\eta_s=\chi_i^{-1}$ のみで、各重みは高々一次元となる。

この三方向は全て同時に実現する。位置 $n$ の $x_lx_k-q_{kl}x_kx_l$ に
$\sum_{s=0}^2 C_s\eta_s^n x_i^s x_j^{2-s}$ を加える（固定順序 $i,j$ を使う）。
新旧各文字を通る全三文字重なりを取る。補正が現れると $k,l$ がともに消えるので
同じ書換えを再使用しない。従って上の線型条件だけで差は零である。
$k,l$ の個数を先に、次に反転数を下げて停止するため、全係数で順序語が基底になる。
量子窓への平坦退化またはこの二次PBW表示と量子外積への双対の随伴次数により標準正則である。
$\eta_s$ の比は連続する $s$ で一定で、四周期の各係数倍率は三指標への共通トーラス作用で実現できる。
従って一位置移動は各係数で同型になる。
\end{proof}

\begin{proof}[候補の尽くし方と同型]
量子テンソルの全24個の置換係数は非零である。この台の対角基底変更に関する重みの重心は
各頂点の四座標で均等なので、その重心は重み凸包の相対内部にある。
よって複素数上で四つの正対角行列により各一文字切断の四行のノルムを同じにできる。
具体的には跡零対数倍率上のノルム平方は有限個の正係数指数関数の和で、
有効な部分では凸包の内部条件から無限遠で発散し最小値を持つ。
異なる行の台は交わらないので全Gram行列がスカラーとなり、量子点はpolystableである。
代数的な閉軌道性は任意の標数零代数閉体に移る。
面値1が高々一つの範囲では四座標点が点スキームから内在的に復元され、
連結安定化群は共通対角トーラスである。異なる閉軌道は同じ商ファイバーを持たない。

三係数族のFrobenius標準四移動は普通の四反復である。
実際、純側 $i,j$ に重み1、$k,l$ に重み0を付けると、係数は全て重み2にできる。
PBW最高次の基底で規格化したNakayama行列は係数に正則で、各成分の重みの絶対値は高々1である。
従って基点から変化できない。最高次テンソルを量子体積の一つの係数で規格化すれば同様に、
全語の重みは0から4、体積語は重み2であるから補正は三係数の一次までである。
ゆえにテンソル族は基点を通るアフィン三平面となる。三方向は負添字が二つで
文字変更による境界がなく、異なる指標であるから軌道接と交わらない。

前節と同じLuna切片を用いるが、今回は周囲の滑らかさを仮定しない。
量子点の有限置換部分もこの三方向の空間を保つ。実際、反転した順序語の並べ替えでは
純 $i,j$ に補正がなく、$kl$ 以外の関係も変わらない。従って安定化群不変な軌道接の補空間にこの平面を含めて、
アフィン線型横断面と階数スキームの交叉で切片を選べる。射影倍率は体積係数の不変な局所規格化で除く。
切片の完備局所環の位相的生成元は、余接空間のトーラス固有基底を持ち上げて取れる。
上の接計算により重み零のもの以外は、単一重複の場合の三平面の座標だけで足りる。
その三座標は平面上で $C_s$ と一致する線型射影を選べる。重み零の生成元は
この平面に制限すると定数零である。純側二文字だけに同じ正の余重みを与えれば三方向とも正だからである。

切片の基点ファイバーの任意の点 $y$ は一径数部分群で原点へ退化できる。
三座標を $C_s(y)$ とし、同じ三値の平面上の点を取る。
前者が原点へ流れる余指標ではその非零座標は全て正なので、後者も同じ余指標で流れる。
二つの弧では重み零生成元が定数零、他の三座標が同じ単項式となるため、
全完備局所関数が一致する。従って二つの弧は一致し、元の点も同じである。
ここで切片の飽和エタール近傍は基点ファイバー全体を含む。横断平面上の全点も原点への
一径数退化を持つため同じ近傍に入る。四面値が相異なる場合は三座標もなく、同じ議論で基点だけが残る。

三指標の整数関係は $\alpha_0+\alpha_2=2\alpha_1$ の整数倍だけであり、
二つ以下では独立である。従ってトーラス商の同型条件は所定の台と単項式比である。
単一重複の場合、偶置換の安定化元は恒等だけである。奇置換では重複値が自分の逆となる必要があり、
$c\ne1$ より $c=-1$ となる。残る二値は相異なる逆数なので可能な置換は $(ij)$ だけである。
これを位置0の純置換とし、位置1で追加倍率 $d_i,d_j$ を使うと
$d_i/d_j=q_{ij}^2$ である。混合単項式の並べ替えの倍率 $q_{ij}$ と合わせれば、
単項式比に掛かる係数 $d_i/(d_jq_{ij}^2)$ は1となる。従ってこの場合も比の値は保たれる。
\path{computations/verify_quantum_face_resonance.py} は面の符号と置換条件を十二乗根格子でも検算する。
\end{proof}
